\documentclass[conference]{IEEEtran}
\usepackage{cite}
\usepackage{amsmath,amssymb,amsfonts}
\usepackage{graphicx}
\usepackage{textcomp}
\usepackage{xcolor}
\usepackage{hyperref}
\usepackage{booktabs}
\usepackage{float}
\usepackage{balance}

\begin{document}

\title{Real-Time Traffic Sign and Signal Detection for Live Video Streams Using YOLOv8}

\author{
\IEEEauthorblockN{Longric Tran, Sheel Patel, Hareesh Suresh}
\IEEEauthorblockA{Department of Computer Science\\
Toronto Metropolitan University\\
Toronto, Ontario, Canada\\
longric.tran@torontomu.ca, s60patel@torontomu.ca, hsuresh@torontomu.ca}
}

\maketitle

%==============================================================================
% ABSTRACT
%==============================================================================
\begin{abstract}
Traffic sign and traffic light recognition is a core perception capability for driver-assistance and autonomous driving systems, but practical deployment requires both reliable accuracy and low latency. In this work, we build a real-time detector by fine-tuning an Ultralytics YOLOv8 model on a 15-class traffic sign dataset containing speed-limit signs (10--120 km/h), a stop sign, and red/green traffic lights. We present the complete pipeline: dataset configuration in YOLO format, transfer-learning based fine-tuning, offline evaluation using standard detection metrics (precision, recall, mAP), and deployment on live webcam streams using OpenCV. We also discuss common failure cases such as far-range small objects, digit ambiguity among speed limits, and domain shift when signs are presented on phone screens. The result is a lightweight prototype that demonstrates the speed--accuracy tradeoff of YOLOv8 and serves as a practical baseline for extending toward more realistic driving scenarios.

\vspace{0.4em}
\noindent\textbf{Source Code:} \url{https://github.com/longrict/cps843final/}
\end{abstract}

%==============================================================================
% INTRODUCTION
%==============================================================================
\section{Introduction}
Traffic signs and regulatory signals communicate critical rules (speed limits, stopping requirements, right-of-way) and safety warnings in a fast and largely language-independent manner. As Advanced Driver Assistance Systems (ADAS) and autonomous vehicles become more prevalent, automatic sign detection becomes increasingly important for safe decision-making: a vehicle must correctly interpret its driving context to enforce speed compliance, handle intersections, and respond to traffic controls.

A deployable detector must satisfy three requirements simultaneously: high \textbf{accuracy} (tight localization and correct class), \textbf{robustness} (working under different illumination, weather, and occlusion), and \textbf{low latency} (real-time responsiveness). Traditional pipelines relied on hand-crafted features (color thresholding, edge detection, HOG/SIFT descriptors) combined with classical classifiers. While effective under controlled conditions, these approaches often degrade under domain shift such as shadows, night-time lighting, motion blur, and camera noise. Deep learning approaches based on Convolutional Neural Networks (CNNs) learn features directly from data and have become the dominant strategy for robust detection.

Object detection differs from classification: a detector outputs both a category and a bounding box for each object instance. One-stage detectors such as the You Only Look Once (YOLO) family predict boxes and classes in a single forward pass, offering strong speed advantages over two-stage methods that rely on region proposals. In this project, we fine-tune a pre-trained YOLOv8 model on a traffic sign dataset with 15 labels and deploy the trained weights in a webcam-based real-time pipeline.

Our main contributions are:
\begin{itemize}
    \item An end-to-end traffic sign detection pipeline using YOLOv8, including dataset configuration, fine-tuning, evaluation, and webcam deployment.
    \item A detailed technical description of YOLOv8 components (backbone/neck/head), training objective, and inference post-processing.
    \item A structured evaluation plan combining offline mAP/precision/recall with real-time FPS/latency measurements and qualitative error analysis.
\end{itemize}

%==============================================================================
% RELATED WORK
%==============================================================================
\section{Related Work}
The original YOLO model unified detection into a single network, enabling real-time object detection by avoiding computationally expensive region proposal stages \cite{YOLO_original}. Modern YOLO implementations follow a backbone--neck--head structure in which a convolutional backbone extracts features, a neck aggregates multi-scale information, and a head produces final class and box predictions \cite{viso_yolov8}.

Two widely used alternatives to YOLO-style one-stage detectors are Faster R-CNN and SSD. Faster R-CNN is a two-stage detector: it generates region proposals and then classifies/refines them, often achieving strong accuracy but at higher computational cost \cite{faster_rcnn}. SSD is a one-stage detector that predicts detections at multiple feature map scales and is typically faster than two-stage methods while maintaining competitive accuracy \cite{SSD}. For traffic sign detection, the primary challenges include small-object detection at far distances, visually similar classes (speed-limit digits), and domain shift between curated datasets and real-world driving video. Multi-scale feature fusion such as Feature Pyramid Networks (FPNs) is commonly used to improve small-object performance \cite{FPN}, and transfer learning is widely adopted to adapt generic pre-trained models to task-specific label spaces.

%==============================================================================
% TECHNICAL PART
%==============================================================================
\section{Technical Approach}

\subsection{YOLO Detection Overview}
YOLO is a one-stage detector: the model predicts bounding boxes and class probabilities directly from the input image in a single pass. Candidate predictions are filtered using a confidence threshold and then de-duplicated using Non-Maximum Suppression (NMS). This design reduces inference overhead and is therefore well-suited for streaming video.

Ultralytics describes YOLOv8 as a modular architecture composed of a \textbf{backbone}, \textbf{neck}, and \textbf{head} \cite{viso_yolov8}. Intuitively, the backbone extracts features, the neck fuses information at different scales, and the head produces final detections.

\subsubsection{Backbone}
The backbone is the primary feature extraction network. It converts a raw image into a hierarchy of feature maps representing edges, textures, and higher-level shapes \cite{backbone}. Ultralytics YOLO models employ CSP-style designs that improve efficiency and gradient flow by splitting feature maps, processing only part through dense blocks, and concatenating results \cite{CSPNet}. This reduces redundant computation and supports faster inference, which is important for real-time deployment.

\subsubsection{Neck}
The neck aggregates features from multiple backbone stages so that both small and large objects can be detected accurately \cite{backbone}. For traffic signs, scale variation is large: signs may occupy only a small fraction of the image when far away. A Feature Pyramid Network (FPN) addresses this by combining high-resolution spatial detail from shallow layers with semantic information from deeper layers using a top-down pathway and lateral connections \cite{FPN}.

\subsubsection{Head}
The detection head makes the final predictions: object class probabilities and bounding box coordinates \cite{head}. Many modern detectors use decoupled branches, where one branch focuses on classification and another on regression. The regression branch aims to maximize overlap between predicted and ground-truth boxes, often measured by Intersection over Union (IoU), while classification aims to maximize the probability assigned to the correct label.

\subsection{Training Objective and IoU}
Let $b$ be a predicted bounding box and $b^\ast$ a ground-truth box. IoU measures overlap:
\begin{equation}
\mathrm{IoU}(b, b^\ast) = \frac{|b \cap b^\ast|}{|b \cup b^\ast|}.
\end{equation}
Training seeks to produce boxes with high IoU and correct class probabilities. While Ultralytics abstracts the exact loss components, the objective can be understood as a weighted combination of a localization loss (related to IoU) and a classification loss. Inference keeps predictions above a confidence threshold and uses NMS to remove duplicates.

\subsection{Non-Maximum Suppression}
Detectors often output multiple overlapping boxes for the same physical sign. NMS removes duplicates by keeping the highest-confidence detection and suppressing other detections that overlap beyond an IoU threshold $\tau$. For detections sorted by confidence, for each detection $d_i$ we remove any $d_j$ with $\mathrm{IoU}(d_i, d_j) > \tau$. Smaller $\tau$ is more aggressive (fewer duplicates but higher risk of suppressing close objects), while larger $\tau$ preserves more boxes but can increase duplicates.

\subsection{Dataset Selection and Label Space}
We use the Traffic Signs Detection dataset by P. K. Darabi \cite{kaggle_dataset}. The dataset is formatted for YOLO: each image has a corresponding text file with one line per object:
\begin{verbatim}
[class_id] [center_x] [center_y] [width] [height]
\end{verbatim}
Coordinates are normalized to $[0,1]$. Images are 416$\times$416 RGB. The dataset contains 15 classes spanning speed limits, traffic lights, and a stop sign (Table~\ref{tab:categories}). We use 3,530 images for training and 801 for validation, with a separate held-out test split for final evaluation \cite{kaggle_dataset}.

\begin{table}[H]
\centering
\caption{Traffic Sign Classes in Dataset}
\label{tab:categories}
\begin{tabular}{ll}
\toprule
\textbf{Category} & \textbf{Classes} \\
\midrule
Traffic Lights & Green Light, Red Light \\
Speed Limits & 10--120 km/h (12 classes) \\
Regulatory & Stop \\
\bottomrule
\end{tabular}
\end{table}

\subsection{Baseline and Need for Fine-Tuning}
We tested the COCO-pretrained YOLOv8 weights on traffic sign images and observed either no detections or detections in unrelated categories. This behavior is expected because COCO contains general object categories and does not include speed-limit sign classes. Fine-tuning adapts the detector to the task-specific label space \cite{kaggle_dataset}.

\subsection{System Pipeline}
Figure~\ref{fig:pipeline} summarizes the end-to-end workflow. The same trained weights are used for both offline evaluation and webcam deployment, ensuring that reported metrics correspond to the deployed model.
\begin{figure}[H]
\centering
\includegraphics[width=\linewidth]{figures/pipeline_placeholder.png}
\caption{End-to-end pipeline: dataset configuration $\rightarrow$ fine-tuning $\rightarrow$ evaluation on held-out data $\rightarrow$ live webcam inference with OpenCV.}
\label{fig:pipeline}
\end{figure}

\subsection{Data Augmentation and Generalization}
Deployment video differs from curated datasets: signs appear at different scales, under shadows, with motion blur, and may be partially occluded. To improve robustness, YOLO training typically uses stochastic augmentations such as random scaling/translation, color jitter (HSV shifts), and mosaic-style composition that combines multiple images into one training sample. These augmentations help the model generalize by exposing it to a larger effective training distribution and reducing overfitting to dataset-specific backgrounds. In our implementation we rely on the default Ultralytics augmentation recipe, which provides a strong baseline without extensive manual tuning.

\subsection{Training Procedure and Configuration}
Training was conducted using the Ultralytics YOLO framework built on PyTorch. We choose a lightweight YOLOv8 variant to prioritize low latency in real-time inference. Unless otherwise noted, optimizer, learning rate schedule, and default augmentations follow the framework recipe. Core hyperparameters are shown in Table~\ref{tab:trainingparameters}.

\begin{table}[H]
\centering
\caption{Training Hyperparameters}
\label{tab:trainingparameters}
\begin{tabular}{ll}
\toprule
\textbf{Parameter} & \textbf{Value} \\
\midrule
Base Model & YOLOv8 (lightweight variant) \\
Epochs & 50 \\
Image Size & 416 \\
Batch Size & 16 \\
Confidence Threshold (infer) & 0.25 \\
IoU Threshold (NMS) & 0.45 \\
\bottomrule
\end{tabular}
\end{table}

\subsection{Implementation Details}
The dataset is specified through a YAML configuration listing the train/validation/test paths, the number of classes, and the class-name mapping. Labels are stored in parallel \texttt{labels/} directories with the same file stem as the corresponding images.

A typical fine-tuning command (conceptually) is:
\begin{verbatim}
yolo detect train model=yolov8n.pt data=data.yaml
     imgsz=416 epochs=50 batch=16
\end{verbatim}
After training, the best checkpoint (selected by validation performance) is used for evaluation and deployment. Offline evaluation uses the built-in validation routine, which computes precision, recall, and mAP with IoU-based matching. For reproducible reporting, we log (i) the exact checkpoint name, (ii) confidence/IoU thresholds used during deployment, and (iii) hardware characteristics since these directly affect FPS.

\subsection{Real-Time Detection via OpenCV}
Deployment performs frame-by-frame inference on a live webcam stream using OpenCV:
\begin{enumerate}
    \item Capture a frame from the camera.
    \item Run YOLOv8 inference to produce candidate boxes and class scores.
    \item Filter by confidence and apply NMS.
    \item Render bounding boxes and labels on the frame and display to the user.
\end{enumerate}
For real-time usability, two aspects matter: \textbf{throughput} (frames per second, FPS) and \textbf{end-to-end latency}. We measure FPS over a window of $N$ frames:
\begin{equation}
\mathrm{FPS} = \frac{N}{t_{\mathrm{end}} - t_{\mathrm{start}}}.
\end{equation}
Latency is measured by timing the full loop (capture $\rightarrow$ inference $\rightarrow$ draw $\rightarrow$ display).

\subsection{Deployment Considerations}
In streaming systems, speed can be limited by more than model inference. Common bottlenecks include frame capture, image format conversion (BGR$\leftrightarrow$RGB), and drawing overlays. To maximize throughput, it is useful to avoid unnecessary copies and keep visualization lightweight. A further improvement is to decouple capture and inference using separate threads so inference always runs on the newest frame. Our prototype uses a simple single-thread loop for clarity, but the same trained weights can be integrated into more optimized pipelines.

%==============================================================================
% EXPERIMENTS
%==============================================================================
\section{Experiments}

\subsection{Experimental Setup}
All experiments were implemented in Python using the Ultralytics YOLO library and OpenCV. Training used the dataset split described earlier and ran for 50 epochs with 416$\times$416 inputs. Evaluation was performed on the held-out validation split and/or a separate test split. For deployment testing, we used a laptop webcam and evaluated responsiveness under indoor and outdoor lighting.

\subsection{Evaluation Metrics}
We report standard object detection metrics:
\begin{itemize}
    \item \textbf{Precision:} $\mathrm{Precision} = \frac{TP}{TP + FP}$
    \item \textbf{Recall:} $\mathrm{Recall} = \frac{TP}{TP + FN}$
    \item \textbf{Average Precision (AP):} area under the precision--recall curve for a class at a chosen IoU threshold.
    \item \textbf{mAP:} mean AP across classes, reported as mAP@0.5 and mAP@0.5:0.95.
\end{itemize}
For real-time performance, we report FPS and per-frame latency.

\subsection{Quantitative Results}
Table~\ref{tab:results} is a compact reporting format for offline detection performance. The values should be copied directly from the Ultralytics validation output to avoid transcription error. Reporting both mAP@0.5 and mAP@0.5:0.95 distinguishes between loose overlap performance and stricter localization quality.
\begin{table}[H]
\centering
\caption{Detection Performance on Held-Out Data (fill with values from logs)}
\label{tab:results}
\begin{tabular}{lcccc}
\toprule
\textbf{Model} & \textbf{Prec.} & \textbf{Rec.} & \textbf{mAP@0.5} & \textbf{mAP@0.5:0.95} \\
\midrule
YOLOv8 (fine-tuned) & -- & -- & -- & -- \\
YOLOv8 (COCO-pretrained, no FT) & $\approx$0 & $\approx$0 & $\approx$0 & $\approx$0 \\
\bottomrule
\end{tabular}
\end{table}

Beyond overall mAP, class-wise AP values provide insight into which categories are difficult. Stop signs and traffic lights are visually distinctive, while speed-limit categories are highly similar and can be confused under blur or glare. A confusion matrix over the 15 classes can further reveal which speeds are most frequently mistaken for each other.


\subsection{Per-Class and Category-Level Analysis}
Because the dataset contains a mixture of visually distinctive classes (stop sign, traffic lights) and visually similar classes (speed limits), it is informative to report performance aggregated by category. Table~\ref{tab:category} provides a template for category-level AP. A category view often highlights whether errors are mainly due to digit ambiguity (speed limits) or due to localization and small-object challenges (traffic lights at distance).

\begin{table}[H]
\centering
\caption{Category-Level Performance (fill with measured AP)}
\label{tab:category}
\begin{tabular}{lcc}
\toprule
\textbf{Category} & \textbf{AP@0.5} & \textbf{AP@0.5:0.95} \\
\midrule
Traffic Lights & -- & -- \\
Speed Limits & -- & -- \\
Stop Sign & -- & -- \\
\bottomrule
\end{tabular}
\end{table}

For speed limits, class-wise AP can identify which digits are frequently confused (e.g., 80 vs.\ 90 under blur). When available, reporting a confusion matrix and including example misclassifications improves interpretability beyond a single mAP number.

\subsection{Robustness Stress Tests}
To evaluate robustness in conditions closer to deployment, we recommend stress tests that vary one nuisance factor at a time:
\begin{itemize}
    \item \textbf{Distance/scale:} record scenes where signs occupy progressively smaller regions of the frame.
    \item \textbf{Lighting:} test under bright sunlight, shade, and indoor lighting to expose over/under-exposure failures.
    \item \textbf{Motion blur:} move the camera to simulate driving-like motion and evaluate stability of detections.
    \item \textbf{Occlusion:} test with partial blocking of the sign (e.g., leaves, poles) to assess resilience.
\end{itemize}
For each condition, record the number of missed detections and the confidence distribution; even qualitative summaries can be grounded by counting representative clips.

\begin{figure}[H]
\centering
\includegraphics[width=\linewidth]{figures/confusion_matrix_placeholder.png}
\caption{Confusion matrix placeholder (replace with your computed confusion matrix for the 15 classes).}
\label{fig:conf}
\end{figure}


\subsection{Training Curves}
Training curves help interpret optimization behavior. Figure~\ref{fig:curves} is a placeholder for loss and mAP curves generated during training. A decreasing training loss combined with increasing validation mAP suggests successful fine-tuning. If training loss decreases while validation mAP stagnates or decreases, this indicates overfitting or insufficient augmentation.
\begin{figure}[H]
\centering
\includegraphics[width=\linewidth]{figures/training_curves_placeholder.png}
\caption{Example training curves (replace with your generated plots): training/validation loss and mAP over epochs.}
\label{fig:curves}
\end{figure}

\subsection{Qualitative Results and Error Analysis}
Qualitative inspection is essential to understand real-world behavior. In our tests, the fine-tuned model detects stop signs and traffic lights reliably when the sign occupies a reasonable fraction of the frame and lighting is moderate. We observed several recurring error modes:
\begin{itemize}
    \item \textbf{Small objects:} distant signs may be only a few pixels wide, leading to missed detections or low-confidence boxes.
    \item \textbf{Digit ambiguity:} speed-limit classes are visually similar; under motion blur or downsampling the digits may not be distinct enough, producing confusion (e.g., 60 vs.\ 80).
    \item \textbf{Occlusion and truncation:} partial visibility can disrupt digit recognition even if the sign shape is detected.
    \item \textbf{Screen-based testing:} signs displayed on a phone screen were less consistently detected, likely due to glare and different color statistics.
\end{itemize}
\begin{figure}[H]
\centering
\includegraphics[width=\linewidth]{figures/qualitative_placeholder.png}
\caption{Qualitative examples (replace with screenshots): correct detections and representative failure cases.}
\label{fig:qual}
\end{figure}

\begin{figure}[H]
\centering
\includegraphics[width=\linewidth]{figures/realtime_demo_placeholder.png}
\caption{Real-time webcam demo screenshot (replace with an image from your run).}
\label{fig:demo}
\end{figure}

\subsection{Real-Time Performance}
The webcam inference pipeline ran smoothly on a standard laptop. When testing with nearby printed signs, detections updated continuously with minimal visible lag, enabling interactive use. Table~\ref{tab:rt} provides a reporting format for measured throughput and latency. If a GPU is available, reporting results for both CPU and GPU highlights deployment scaling with hardware.

\subsection{Runtime Profiling}
To interpret FPS measurements, it is useful to break down the end-to-end loop time into components: frame capture, preprocessing/resizing, model inference, and rendering. A simple timing study can report average milliseconds per stage over $N$ frames. This helps identify whether the bottleneck is inference or overhead in capture/visualization. When inference dominates, improvements come from using a smaller model, lowering input resolution, enabling GPU acceleration, or batching (not usually suitable for real-time). When capture/rendering dominates, improvements come from avoiding expensive conversions and reducing drawing operations.

\begin{table}[H]
\centering
\caption{Real-Time Performance (fill with measured values)}
\label{tab:rt}
\begin{tabular}{lcc}
\toprule
\textbf{Platform} & \textbf{FPS} & \textbf{Latency (ms)} \\
\midrule
Laptop CPU (webcam) & -- & -- \\
Laptop GPU (optional) & -- & -- \\
\bottomrule
\end{tabular}
\end{table}

\subsection{Ablation and Sensitivity Analysis}
To understand which design choices matter most, we recommend ablations that change one factor while holding others constant. Even when absolute numbers vary by hardware and training noise, ablations reveal qualitative trends: higher input resolution can improve far-range detection but reduces FPS; larger models often improve mAP but require more compute; and confidence/NMS thresholds shift the operating point without changing weights. Table~\ref{tab:ablation} is a template for recording these comparisons.
\begin{table}[H]
\centering
\caption{Ablation Templates (fill with measured results)}
\label{tab:ablation}
\begin{tabular}{lcccc}
\toprule
\textbf{Setting} & \textbf{Prec.} & \textbf{Rec.} & \textbf{mAP@0.5} & \textbf{FPS} \\
\midrule
Base (416, 50 ep, light model) & -- & -- & -- & -- \\
Higher res (640, 50 ep) & -- & -- & -- & -- \\
Fewer epochs (416, 25 ep) & -- & -- & -- & -- \\
Larger model (e.g., ``small'') & -- & -- & -- & -- \\
\bottomrule
\end{tabular}
\end{table}

\subsection{Discussion: Accuracy--Latency Tradeoffs}
For ADAS-style applications, the best model is not necessarily the one with the highest mAP; it is the one that meets an end-to-end latency budget while maintaining acceptable detection quality. In our setting, traffic signs can be small and visually similar, which pushes toward higher resolution and stronger backbones, while real-time video pushes toward smaller models. YOLOv8 provides practical knobs (model scale, input size, and inference thresholds) to explore this tradeoff.

Domain shift is also important: the training dataset contains relatively clean, centered signs, while live video includes motion blur, complex backgrounds, and variable illumination. Our qualitative inspection suggests that the most common failures come from insufficient pixel resolution for distant signs and digit ambiguity for speed limits. Both issues can be mitigated by collecting additional training data that matches deployment conditions and by experimenting with higher-resolution training when hardware permits.

\subsection{Limitations and Future Work}
While this project demonstrates a functional prototype, deploying sign detection in safety-critical systems requires rigorous testing, redundant sensing, and calibrated confidence estimates. False negatives (missed stop signs or red lights) are especially concerning in real driving scenarios, and false positives can cause unnecessary braking or driver distraction. Our work should therefore be seen as a course-scale demonstration and a baseline for further engineering rather than a certified safety component.

Several concrete directions can improve performance:
\begin{itemize}
    \item \textbf{More diverse data:} add examples with night lighting, rain, heavy shadows, and partial occlusions.
    \item \textbf{Higher-resolution training:} retrain with larger \texttt{imgsz} and multi-scale augmentation for small-object detection.
    \item \textbf{Reframing speed limits:} detect speed-limit signs and then read digits using OCR, potentially reducing confusion among similar speeds.
    \item \textbf{Temporal smoothing:} incorporate tracking or temporal filtering to reduce flicker across frames.
\end{itemize}

%==============================================================================
% CONCLUSION
%==============================================================================
\section{Conclusion}
We developed a real-time traffic sign and signal detector by fine-tuning a pre-trained YOLOv8 model on a 15-class dataset of speed limits, traffic lights, and stop signs. The resulting system integrates dataset configuration, training, offline evaluation, and webcam inference in a single pipeline using Ultralytics and OpenCV. The main engineering takeaway is that transfer learning is essential: a general-purpose COCO model does not recognize traffic sign categories, but fine-tuning on a task-specific dataset adapts the detector to the desired label space. From a deployment perspective, the key requirement is balancing accuracy with latency; lightweight YOLOv8 variants can remain responsive on commodity hardware while producing meaningful detections for interactive demonstrations. Our qualitative analysis highlights the primary limitations of the current model: distant signs become too small to detect reliably, and speed-limit digits can be ambiguous under blur or glare. These observations point to the most promising improvements: collecting more diverse training data, increasing input resolution when hardware allows, and incorporating temporal smoothing or a dedicated digit-reading stage for speed limits.

\vspace{0.4em}
\noindent\textbf{Team contributions:} Longric Tran led dataset configuration and training orchestration. Sheel Patel implemented the webcam inference pipeline and conducted qualitative testing. Hareesh Suresh prepared the technical write-up, evaluation methodology, and analysis, and coordinated final report integration.

\balance
%==============================================================================
% REFERENCES
%==============================================================================
\begin{thebibliography}{00}

\bibitem{viso_yolov8}
Viso.ai, ``YOLOv8: A Complete Guide,'' 2023. [Online]. Available: \url{https://viso.ai/deep-learning/yolov8-guide/}

\bibitem{YOLO_original}
J. Redmon, S. Divvala, R. Girshick, and A. Farhadi, ``You Only Look Once: Unified, Real-Time Object Detection,'' 2015. [Online]. Available: \url{https://arxiv.org/abs/1506.02640}

\bibitem{faster_rcnn}
S. Ren, K. He, R. Girshick, and J. Sun, ``Faster R-CNN: Towards Real-Time Object Detection with Region Proposal Networks,'' in \emph{Advances in Neural Information Processing Systems (NIPS)}, 2015. [Online]. Available: \url{https://arxiv.org/abs/1506.01497}

\bibitem{SSD}
W. Liu, D. Anguelov, D. Erhan, C. Szegedy, S. Reed, C.-Y. Fu, and A. C. Berg, ``SSD: Single Shot MultiBox Detector,'' in \emph{ECCV}, 2016. [Online]. Available: \url{https://arxiv.org/abs/1512.02325}

\bibitem{backbone}
Ultralytics, ``Backbone.'' [Online]. Available: \url{https://www.ultralytics.com/glossary/backbone}

\bibitem{CSPNet}
C. Wang, H. Liao, I. Yeh, Y. Wu, P. Chen, and J. Hsieh, ``CSPNet: A New Backbone that can Enhance Learning Capability of CNN,'' 2019. [Online]. Available: \url{https://doi.org/10.48550/arXiv.1911.11929}

\bibitem{FPN}
Ultralytics, ``Feature Pyramid Network (FPN).'' [Online]. Available: \url{https://www.ultralytics.com/glossary/feature-pyramid-network-fpn}

\bibitem{head}
Ultralytics, ``Detection Head.'' [Online]. Available: \url{https://www.ultralytics.com/glossary/detection-head}

\bibitem{kaggle_dataset}
P. K. Darabi, ``Traffic Signs Detection Using YOLOv8,'' Kaggle, 2023. [Online]. Available: \url{https://www.kaggle.com/code/pkdarabi/traffic-signs-detection-using-yolov8}

\end{thebibliography}

\end{document}
